\documentclass[11pt]{article}

\usepackage{fullpage}
\usepackage{enumitem}
\usepackage{amsmath,amsthm,amsfonts,amssymb,mathtools,latexsym}
\usepackage{xcolor,graphicx,hyperref,float}
\usepackage{tikz}
\usetikzlibrary{arrows.meta}
\usepackage[ruled,vlined,linesnumbered]{algorithm2e}

\pdfpagewidth 8.5in
\pdfpageheight 11in
\oddsidemargin 0in
\evensidemargin 0in

\providecommand{\tightlist}{\setlength{\itemsep}{0pt}\setlength{\parskip}{0pt}}
\providecommand{\squigrightarrow}{\rightsquigarrow}

\begin{document}
\begin{center}
\Large{\textbf{Homework 8}}
\end{center}

\hypertarget{maximum-clearance-routes}{%
\section{1. Maximum Clearance Routes}\label{maximum-clearance-routes}}

\hypertarget{a-cut-and-cycle-property-of-maximum-spanning-tree.}{%
\subsection{\texorpdfstring{(a) Cut and Cycle Property of \emph{maximum}
spanning
tree.}{(a) Cut and Cycle Property of maximum spanning tree.}}\label{a-cut-and-cycle-property-of-maximum-spanning-tree.}}

\[
\text{Maximum Spanning tree G: a tree that connects all vertices in G }
\text{while maximizing the sum of the edge weights in the tree.}
\]

\begin{enumerate}
\def\labelenumi{\arabic{enumi}.}
\item
  Cut Property: Let \(i\) and \(j\) be the nodes in the two sides of a
  cut. Let \(e\) be the maximum weight across the cut. Then \(e\) is in
  the Maximum Spanning Tree \(G^{*}\)
\item
  Cycle Property: Let \(C\) be a cycle and let \(f\) be the minimum
  weight edge in \(C\). Then, \(f\) is \emph{not} in the Maximum
  Spanning Tree.
\end{enumerate}

\hypertarget{b-maxst-proof}{%
\subsection{(b) MaxST Proof}\label{b-maxst-proof}}

Let \(a\) be another route between \((u,v) \in G\) that has a greater
route weight than the original edge between \((u,v)\), \(o \in G\) and
that it creates a cycle. Then, according to the cycle property of MaxST,
we have to choose \(a\) route to be MaxST. Here, using \(o\) route
contradicts the cycle property. Since \(o\) is in the MaxST, \(o\) must
have the maximum route weight; otherwise there would be contradiction in
the cycle property. Therefore, any graph \(G\) with unique edge
clearance and disntinct route weights, the path given by using edges
from the maximum spanning tree of \(G\) gives the maximum clearance
route of \(G\).

\hypertarget{c-maxst-polynomial-runtime-algorithm}{%
\subsection{(c) MaxST Polynomial runtime
Algorithm}\label{c-maxst-polynomial-runtime-algorithm}}

\hypertarget{pseudocode}{%
\subsubsection{Pseudocode}\label{pseudocode}}

\begin{algorithm}[H]
    \small
    \caption{MaxST($G$)}
    \tcc{$G$: undirected, weighted adj list. Kruskal's Algoritm with edges sorted by decreasing order}
    return ReverseKruskal($G$)\;
\end{algorithm}

\hypertarget{runtime-analysis}{%
\subsubsection{Runtime Analysis}\label{runtime-analysis}}

The above pseudocode runs Kruskal's Algorithm with sorted edges by
decreasing order. Since the sorting edges reversely does not change the
runtime, the total runtime of the above algorithm is same as Kruskal's
Algorithm's runtime: \(O(mlog(m))\).

\hypertarget{proof-of-correctness}{%
\subsubsection{Proof of Correctness}\label{proof-of-correctness}}

For the algorithm to always output the maximum spanning tree,

\begin{enumerate}
\def\labelenumi{\arabic{enumi}.}
\item
  the output tree must correctly have the cut and cycle properties of
  MaxST.

  \begin{itemize}
  \tightlist
  \item
    the output tree must be acyclic.
  \end{itemize}
\item
  all nodes must be connected
\item
  the sum of the edge weights of the tree in the graph \(G\) must be
  greater than or equal to the sum of edge weights of any other spanning
  tree in the graph \(G\).
\end{enumerate}

According to the proof of correctness of Kruskal's algorithm's first
case that checks the existence of cycle, there exists no edge that
creates cycle such that the output tree is always acyclic.

Moreover, for all nodes to be connected, there needs to be \(n-1\) edges
when there exists \(n\) nodes in the tree from the graph \(G\).
Kruskal's Algorithm ensures to stop searching when \(|T| = n-1\) where
\(T\) is the output tree as we are always given the graph \(G\) with
distinct weights for each edge in the graph which ensures no overlaps in
edge weights. Reversing the node alignment for searching does not affect
the stop condition of the algorithm. Therefore, all nodes are always
connected in the output tree.

Lastly, the algorithm greedly selects nodes that does not create cycles
from sorted edges in decreasing order. According to the cycle and cut
properties of MaxST, there cannot be edges with minimum weights in MaxST
and there must be edges with maximum weights in MaxST. Therefore, the
sum of the edge wegiths in the output tree must be greater than or equal
to any other possible spanning tree from the graph \(G\).

Therefore, \texttt{MaxST(G)} always outputs the correct maximum spanning
tree from any arbitrary graph \(G\) that has dintinct edge weights in
each edge of it.

\hypertarget{streaming-edges}{%
\section{Streaming Edges}\label{streaming-edges}}

\hypertarget{a}{%
\subsection{(a)}\label{a}}

Problem: Given \(G, T, (u,v), w(e)\), how is it possible to get a MST
\(T'\) for the larger graph \(G'\) in time \(O(n)\) where \(n\) is the
number of nodes in \(G\).

\hypertarget{facts}{%
\subsubsection{Facts}\label{facts}}

\begin{enumerate}
\def\labelenumi{\arabic{enumi}.}
\item
  \(G = (V,E)\): undirected graph with edge weights \(w\)
\item
  Already computed MinST, \(T\) for \(G\).
\item
  \(G'\) = \(G\) with a new edge \(e = (u,v)\) with weight \(w(e)\)
\item
  \(T'\) is given as a graph with the same node set as \(G\)
\item
  \(O(n)\) is not enough to re-read the original graph \(G\). Have to
  explain why OK to ignore lots of the edges in \(G\).
\end{enumerate}

\hypertarget{proof}{%
\subsubsection{Proof}\label{proof}}

Getting \(G'\) by adding additional edge \(e=(u,v)\) with weight
\(w(e)\) from the original graph \(G\) between any existing nodes in
\(G\), \(u,v\) creates a cycle since MST from any graph is unique.

To check whether we should keep the new edge \((u,v)\) and which edge to
delete iff the new edge \((u,v)\) has the minimum weight wihtin the
cycle \(C\), check all nodes that has edges and that create the cycle
\(C\) from either \(u\) or \(v\) (it does not matter where to check
from).

\begin{itemize}
\item
  check all nodes from either \(u\) or \(v\) and see if the edge weight
  between either \(u\) or \(v\) and any nodes in \(C\) is greater than
  \((u,v)\).

  \begin{itemize}
  \item
    if greater than \((u,v)\): delete the heaviest edge and keep
    \((u,v)\) from \(T\) to create new \(T'\)
  \item
    else: the original \(T\) is \(T'\).
  \end{itemize}
\end{itemize}

Because the worst case of this solution checks all nodes in the graph,
it runs \(O(n)\) in worst case.

\hypertarget{b}{%
\subsection{(b)}\label{b}}

Write a MST algorithm with lazy operations with streams.

\begin{enumerate}
\def\labelenumi{\arabic{enumi}.}
\tightlist
\item
  \(edge_stream\): stream data structure
\end{enumerate}

\begin{itemize}
\item
  \(edge_stream.has_next()\):

  \begin{itemize}
  \item
    return type: \texttt{boolean}
  \item
    checks if there are more edges in the stream
  \end{itemize}
\item
  \(edge_stream.next_edge()\):

  \begin{itemize}
  \item
    return type: \texttt{pair}
  \item
    returns the next edge in the stream
  \end{itemize}
\end{itemize}

\begin{enumerate}
\def\labelenumi{\arabic{enumi}.}
\setcounter{enumi}{1}
\tightlist
\item
  cannot go backwards and re-read previous edges
\end{enumerate}

\begin{itemize}
\tightlist
\item
  store them if we want to re-access
\end{itemize}

\hypertarget{pseudocode-1}{%
\subsubsection{Pseudocode}\label{pseudocode-1}}

\begin{algorithm}[H]
    \small
    \caption{LazyMST($V$, $edge\_stream$)}
    \tcc{overall idea:}
    \tcc{1. compare weights of its available streamed edges}
    \tcc{2. pick the minimum edge between u and vand exclude the edges that cannot be in MST or the ones we already checked}
    \tcc{3. add it to MST}
    $MST \gets$ empty adjacency list\;
    \tcc{first, we need to initialize MST adj list with nodes: space: $O(n)$, time: $O(n)$}
    \ForEach{each vertex $n$ in $V$}{
        $MST[n] \gets$ empty list\;
    }
    \tcc{checking for each edge: $O(m)$}
    \While{$edge\_stream.has\_next()$}{
        $(u,v,w) \gets edge\_stream.next\_edge()$\;
        $path \gets \text{partialDFS}(MST,u,v)$\;
        \tcc{the path $u \squigrightarrow v$ found by DFS. Additional edge weight tracking added}
        \eIf{$path = null$}{
            $MST[u] \gets (v,w)$\;
            $MST[v] \gets (u,w)$\;
        }{
            $edge \gets path[0]$\;
            $max\_weight \gets edge[2]$\;
            \ForEach{$parent, child, weight$ in $path$}{
                \If{$weight > max\_weight$}{
                    $max\_weight \gets weight$\;
                    $edge \gets (parent, child, weight)$\;
                }
            }
            \tcc{this ensures that we do not check for edges that we already checked or}
            \tcc{that cannot be in MST so that we do not have to check it ever in the stream}
            \If{$w < max\_weight$}{
                remove $edge$ from both endpoint lists in $MST$\;
                $MST[u] \gets (v,w)$\;
                $MST[v] \gets (u,w)$\;
            }
        }
    }
    return $MST$\;
\end{algorithm}

\begin{algorithm}[H]
    \small
    \caption{partialDFS($MST$, $u$, $v$)}
    \tcc{Run DFS starting at u in the current forest. Each MST[u] entry is}
    \tcc{(connecting nodes, weight). Return the path as a list of connected edges.}
    \tcc{variation of original DFS; it finds the path from u to v only.}
    $parent\_edge \gets$ hash table\;
    \tcc{we need hash table that keeps track of edges between nodes on top of family relationship}
    $found \gets \text{DFS-visit}(u,v)$\;
    \tcc{boolean value to see if there exists a path}
    \If{$found$}{
        \tcc{if we have path, we need the subgraph S to find the optimal path in LazyMST}
        $path \gets$ empty list\;
        \tcc{insert all the connecting nodes between u and v}
        $con\_nodes \gets v$\;
        \While{$con\_nodes \ne u$}{
            $path.append(parent\_edge[con\_nodes])$\;
            $con\_nodes \gets parent\_edge[con\_nodes][0]$\;
            \tcc{move to the parent endpoint}
        }
        \$\; return path\;
    }
    \tcc{if not found, then null}
    return null\;
\end{algorithm}

\begin{algorithm}[H]
    \small
    \caption{DFS-visit($u$, $v$)}
    \tcc{base case to stop recursion}
    \If{$u = v$}{
        return true\;
    }
    \tcc{recursively search all nodes that from u to v}
    \tcc{we may have connecting nodes between u and v}
    \ForEach{$(con\_node,w)$ in $MST[u]$}{
        \tcc{If the connecting node has not been visited, record its parent edge.}
        \If{$con\_node \notin disc$}{
            $parent\_edge[con\_node] \gets (u,con\_node,w)$\;
            \If{$\text{DFS-visit}(con\_node,v)$}{
                return true\;
            }
        }
    }
    return false\;
\end{algorithm}

\hypertarget{runtimespace-analysis}{%
\subsubsection{Runtime/Space Analysis}\label{runtimespace-analysis}}

\hypertarget{space}{%
\paragraph{Space}\label{space}}

\begin{enumerate}
\def\labelenumi{\arabic{enumi}.}
\item
  \(O(n)\): we use adjacency list \texttt{MST} to store \(n\) nodes
\item
  \(O(n)\): \texttt{parent\_edge} takes \(n\) number of edge entries at
  most in stream. For each node search for minimum weight, it gets
  initialized such that it never exceeds \(O(n)\). Therefore the space
  complexity of this algorithm is \(O(n)\).
\end{enumerate}

\hypertarget{runtime}{%
\paragraph{Runtime}\label{runtime}}

\begin{enumerate}
\def\labelenumi{\arabic{enumi}.}
\item
  \(O(n)\): MST initialization with \(n\) entries
\item
  \(O(m)\): outer loop iterating through \(m\) edges in the stream
\item
  \(O(n)\): \texttt{partialDFS} runs \(O(n)\) times for each iteration
  of outer loop
\item
  \(O(m)\): parallel to the \texttt{partialDFS} recursive call; in worst
  case, we iterate through the entire edge.
\end{enumerate}

\(O(n) + O(m) \cdot (O(n) + O(m) (\text{where n > m})) = O(nm)\)
Therefore, the total runtime of this algorithm is \(O(nm)\).

\hypertarget{proof-of-correctness-1}{%
\subsubsection{proof of correctness}\label{proof-of-correctness-1}}

For this algorithm to always output the correct MST, the MST for each
processed streamed edge must be a minimum spanning forest of the graph
given formed by i edges where i is the number of processed, streamed
edges so far.

\hypertarget{proof-by-induction}{%
\paragraph{Proof by induction}\label{proof-by-induction}}

\hypertarget{base-case}{%
\subparagraph{Base case}\label{base-case}}

if the forest is empty, then it is a minimum spanning forest of the
empty graph

\hypertarget{inductive-step}{%
\subparagraph{Inductive Step}\label{inductive-step}}

\begin{itemize}
\item
  for \(k\)th edge \((u_k, v_k, w_k)\), if \(u_k\) and \(v_k\) are not
  connected yet in MST, connect.
\item
  Otherwise, there will be cycle.
\item
  for any other edges that is on the path from \(u\) to \(v\), if adding
  it to MST creates a cycle, by the cycle property of MST, the edge with
  maximum weight within the comparison cannot be in MST.
\item
  So, we exclude those edges from the stream such that the minimum edge
  weight gets preserved in each edge of the forest.
\item
  So, after the terminating the stream as we went through all possible
  edges in the graph, there will only be edges with minimum edge weights
  left.
\end{itemize}

Therefore, the algorithm always correctly computes MST.

\end{document}
